\section{Research use and interpretation}

Research use should begin with a mechanism-specific question and an explicitly defined observation unit. Table~\ref{tab:research_opportunities} connects the eight branches to research questions, available observations and additional evidence requirements. These mappings identify opportunities for reuse rather than report tests of the proposed hypotheses. Analyses should retain implementation version, actor role, asset, economic stage and observation window, and establish comparability for the particular measure being studied. A shared schema facilitates this work but does not, by itself, identify causal effects or make different native activities interchangeable.

\begin{table}[!htbp]
\centering\tabstyle[registrytint]
\caption{Research opportunities supported by the eight branches. Rows identify potential uses and evidence needed for stronger claims; they are not findings or tests of mechanism-design guarantees.}
\label{tab:research_opportunities}
\begin{tabularx}{\linewidth}{@{}L{0.13\linewidth} L{0.25\linewidth} L{0.26\linewidth} Y@{}}
\toprule\tablehead
\textbf{Branch / literature} & \textbf{Question} & \textbf{Available observations} & \textbf{Further evidence needed} \\
\midrule
UMA~\cite{cong2025oracleprimer,uma2026} & How do challenges and settlement timing vary with incentive commitments? & Bonds, configured rewards, liveness and request lifecycles. & Monitoring opportunities, costs, factual labels and an identification design for incentive effects. \\
Chainlink~\cite{liu2021firstlook,kaleem2021pythia,chainlinkstaking2026} & How do participation and reward realization relate to service provision? & Stake operations, reward claims and selected ETH/USD updates. & Node--service mappings, historical obligations and external references. \\
Tellor~\cite{cong2025oracleprimer,tellorDisputes2026} & How is accountability executed in observed disputes? & Reports, disputes, votes, payments and restrictions. & Broader comparable disputes, challenge opportunities and independent factual outcomes. \\
Flare FTSOv2 \mbox{\cite{cong2025oracleprimer,flareFTSO2026}} & How are reward entitlements formed and realized? & Provider conditions, entitlements, claims and epoch signing policies. & Eligibility--claim matching; historical runtime binding and entity-level ownership for control analysis; external prices for accuracy. \\
Pyth~\cite{cong2025oracleprimer,pythOIS2026} & How do participation and economic realization vary across OIS regimes? & Retained publisher--epoch states, staking instructions and transfers. & Effective rule-change times, comparable windows and treatment of concurrent changes. \\
Chronicle~\cite{eskandari2021sok,chronicleDocs2026} & Which reporting and challenge pathways are observable? & Registered reports, challenges and challenge rewards. & Historical signer sets, thresholds and factual references; full compensation remains unobserved. \\
RedStone Push~\cite{eskandari2021sok,redstoneDocs2026} & How do signed publications translate into Push updates? & Publication and update records for registered adapters and feeds. & Historical update settings and external references; AVS economics requires separate acquisition. \\
DIA~\cite{cong2025oracleprimer,diaDocs2026} & How much of a withdrawal is capital recovery versus reward? & Staking withdrawals and principal/reward components. & Deposit--withdrawal matching and costs for position returns; feeder lineage for price-production analysis. \\
\bottomrule
\end{tabularx}
\end{table}


DuckDB, PyArrow and compatible tools can query Parquet tables through column selection and filtering. The repository notebook provides a smaller reconstruction example; the catalogue and dictionaries document table relationships and variables. Monetary values require asset and decimal fields. Event order follows native chronology, not row position. Addresses identify accounts, not their owners.

Economic aggregation must distinguish configured rewards, accruals, entitlements, payments and applied penalties. Returned principal differs from gain, and net profit requires costs. Check related components before summation; keep USDC.e and USDC separate unless a justified conversion rule applies. Missing, unavailable and inapplicable values are not zeros. Retain censored lifecycles as needed.

The UMA task supports modelling subject to retrospective selection and decision-time constraints. Causal analysis requires an identification strategy and further evidence.

Repositories provide data, code, documentation and an issue tracker. Reuse should record input versions and the computational environment, distinguish archived reconstruction from historical reacquisition, and preserve native meanings when extending analyses or reporting corrections.
